% ===================== ALCANCE =====================
\section{Alcance y Limitaciones}
\label{sec:alcance}

\subsection{Alcance del Estudio}

[Definir con precision los elementos incluidos en el alcance del estudio.
Establecer los limites del trabajo: que se analiza, que sistemas/equipos
estan cubiertos, que metodologias se emplean, y que entregables se producen.]

\begin{table}[htbp]
    \centering
    \caption{[Titulo descriptivo de los elementos incluidos en el alcance.]}
    \label{tab:alcance_incluido}
    \setcounter{itemcount}{0}
    \begin{tabularx}{\textwidth}{@{}NX@{}}
        \toprule
        \multicolumn{1}{c}{\textbf{Item}} & \textbf{Descripcion} \\
        \midrule
        & <Elemento 1 incluido en el alcance.> \\
        & <Elemento 2 incluido en el alcance.> \\
        & <Elemento 3 incluido en el alcance.> \\
        & <Elemento 4 incluido en el alcance.> \\
        \bottomrule
    \end{tabularx}
\end{table}

\subsection{Limitaciones del Estudio}

[Declarar explicitamente los elementos excluidos del alcance, las restricciones
bajo las cuales se ejecuto el trabajo (tiempo, informacion disponible, acceso
a planta), y los supuestos que condicionan los resultados.]

\begin{table}[htbp]
    \centering
    \caption{[Titulo descriptivo de los elementos excluidos del alcance.]}
    \label{tab:alcance_excluido}
    \setcounter{itemcount}{0}
    \begin{tabularx}{\textwidth}{@{}NX@{}}
        \toprule
        \multicolumn{1}{c}{\textbf{Item}} & \textbf{Descripcion} \\
        \midrule
        & <Elemento 1 excluido del alcance.> \\
        & <Elemento 2 excluido del alcance.> \\
        & <Elemento 3 excluido del alcance.> \\
        & <Elemento 4 excluido del alcance.> \\
        \bottomrule
    \end{tabularx}
\end{table}
